% ---------
% --------
%!TEX TS-program = pdflatex
%!TEX encoding = UTF-8 Unicode
%!TEX spellcheck = English (Aspell)
\documentclass[11pt]{article}
\usepackage{lucy,handout,graphicx}
\usepackage[utf8]{inputenc}		% Allow some non-ASCII Unicode in
\usepackage{tikz}

\def\Cdot{\mathbin{\text{\normalfont \textbullet}}}
\def\Sym#1{\texttt{\color{BrickRed}#1}}
\def\BlueSym#1{\textbf{\texttt{\color{RoyalBlue}#1}}}
\allowdisplaybreaks

\newtheorem{theorem}{Claim}

% =========================================================
%   Define common stuff for solution headers
% =========================================================
\Class{CSCI 332}
\Semester{Fall 2025}

% =========================================================
\begin{document}

% ---------------------------------------------------------


\headers{CSCI 332}{Homework 2 Solutions}{Fall 2026}

\begin{enumerate}

\item  In class, we said that the exact running time of an algorithm is the number of primitive operations it
takes on an input of size $n$. (Usually, we disregard this precise running time and instead give some simpler function
such that the precise function is big O if the other function.)
The course textbook (Algorithms by Jeff Erickson) suggests defining the precise running time in a different way.
What is that different way?

\begin{Solution}
    The textbook suggests identifying a specific milestone in the algorithm so that the number of times that milestone
    is reached will be proportional to the overall runtime, regardless of how the algorithm is actually implemented.
\end{Solution}

\item

Consider the following basic problem. You're given an array $A$ consisting of
$n$ integers $A[1], A[2], \ldots, A[n]$. You'd like to output a two-dimensional
$n$-by-$n$ array $B$ in which $B[i,j]$ (for $i < j$) contains the sum of array
entries $A[i]$ through $A[j]$--that is, the sum $A[i] + A[i+1] + \cdots + A[j]$.
(The value of array entry $B[i,j]$ is left unspecified whenever $i \geq j$, so
it doesn't matter what is output by those values.)

\item  Give the resulting $B$ for $A=[4, -9, 12, 2]$.

\begin{Solution}
    \begin{align*}
B[1,2] = -5, B[1,3] = 7, B[1,4] = 9\\
B[2,3] = 3, B[2,4] = 5 \\
B[3,4] = 14
    \end{align*}



\end{Solution}

\begin{algo}
   For $i = 1, 2, \ldots, n$\+ \\
   For $j = i + 1, i + 2, \ldots, n$\+ \\
   Add up array entries $A[i]$ through $A[j]$\\
   Store the result in $B[i,j]$\- \\
   Endfor\- \\
   Endfor
\end{algo}

\begin{enumerate}

    
\item 
Give a function $f$ that counts the number of primitive operations performed by
the algorithm on an input of size $n$. This should be an exact count, not an
asymptotic one!

\begin{Solution}
    Let's try to just count the number
    of additions that happen.  We won't account for other assignment and
    retrieval operations.
    
    As a warm up, let's ask ourselves exactly how many times
    the inner loop runs.
    The outer loop runs exactly $n$ times, for $i = 1, 2, \ldots, n$. The inner loop runs
    $n - i$ times for each value of $i$, since $j$ takes values from $i + 1$
    through $n$. So the total number of times the inner loop runs is
    \[
        (n - 1) + (n - 2) + (n - 3) + \cdots + 1 = \sum_{i=1}^{n-1} i = \frac{(n-1)n}{2}
    \]

    by application of Euler's formula, which says that $\sum_{i=1}^{n} i = \frac{n(n+1)}{2}$.

    Notice that we could have also expressed the number of times the inner loop runs as a double sum:

    \[
    \sum_{i=1}^{n} \sum_{j=i+1}^{n} 1,
    \]
    meaning that $\sum_{i=1}^{n} \sum_{j=i+1}^{n} 1 = \frac{(n-1)n}{2}$, which we will use later.

    Now that we're warmed up, let's count the number of additions.

    Inside the inner loop, we need to add up the entries from $A[i]$ through
    $A[j]$. This requires $(j - i + 1)$ additions. So the total number of
    additions performed by the algorithm is 
    \[
        \sum_{i=1}^{n} \sum_{j=i+1}^{n} (j - i + 1).
    \]

    It would be fine to just look up a way to simplify this summation, but I provide a way to do so here.
    
    We can simplify this expression as follows:

    \begin{align*}
        \sum_{i=1}^{n} \sum_{j=i+1}^{n} (j - i + 1) &= 
        \sum_{i=1}^{n} \sum_{j=i+1}^{n} j - \sum_{i=1}^{n} \sum_{j=i+1}^{n} i + \sum_{i=1}^{n} \sum_{j=i+1}^{n} 1 \\
   \end{align*}

Let's tackle the easiest term first: the last term is just the number of times
the inner loop runs, which we computed above to be $\frac{(n-1)n}{2}$.

    The middle term is a bit trickier. Notice that for each fixed $i$, the inner
    sum $\sum_{j=i+1}^{n} i$ is just $i$ added to itself $(n - i)$ times, or
    $i(n - i)$. So we have
    \[
         \sum_{i=1}^{n} \sum_{j=i+1}^{n} i = \sum_{i=1}^{n} i(n - i) = \sum_{i=1}^{n} ni - i^2 = n\sum_{i=1}^{n} i - \sum_{i=1}^{n} i^2.
    \]
    We can apply Euler's formula to the first term and the sum of squares
    formula ($\sum_{i=1}^{n} i^2 = \frac{n(n+1)(2n+1)}{6}$) to the second term
    to get
    \[
         n\sum_{i=1}^{n} i - \sum_{i=1}^{n} i^2 = n \cdot \frac{n(n+1)}{2} - \frac{n(n+1)(2n+1)}{6}.
    \]
    
    The first term is similar but a bit more complicated.
    For each fixed $i$, the inner sum $\sum_{j=i+1}^{n} j$ is just the sum of
    the integers from $(i+1)$ through $n$. Notice that each of these individuals sums
    can be computed as the sum of $i$ to $n$ with the sum from $i$ to $j$ subtracted off. 
    This is equal to $\sum_{j=1}^{n} j -
    \sum_{j=1}^{i} j$, which we can evaluate by applying Euler's formula to each term. So we have
    \begin{align*}
         \sum_{i=1}^{n} \sum_{j=i+1}^{n} j &= 
         \sum_{i=1}^{n} \left( \sum_{j=1}^{n} j - \sum_{j=1}^{i} j \right) \\ 
         &= 
         \sum_{i=1}^{n} \left(  \frac{n(n+1)}{2} -  \frac{i(i+1)}{2}  \right) \\ 
         &= 
         \sum_{i=1}^{n} \left(  \frac{n(n+1)}{2} -  \frac{i^2}{2} - \frac{i}{2}  \right) \\ 
         &= 
         n\cdot\frac{n(n+1)}{2} -  \frac{1}{2}\left(  \sum_{i=1}^{n} i^2 +  \sum_{i=1}^{n} i  \right) \\
         &= 
         n\cdot\frac{n(n+1)}{2} -  \frac{1}{2}\left(   \frac{n(n+1)(2n+1)}{6} +  \frac{n(n+1)}{2}  \right) \\
    \end{align*}

    Now that we've written each term as a function of $n$, we could put them all
    together and simplify to get our final answer (but I'm not that patient).  Notice that the final
    answer would be $\Theta(n^3)$.

\end{Solution}

\item 
Choose a function $g$ that has the following properties:
    \begin{itemize}
        \item The $f(n)$ you gave above is $O(g(n)$
        \item The $f(n)$ you gave above is $\Omega(g(n))$
    \end{itemize}
\begin{Solution}
    As noted above, our $f(n)$ is clearly both $O(n^3)$ and $\Omega(n^3)$ since it is
    a polynomial of degree three. So $g(n) = n^3$.
\end{Solution}

\item

Although the algorithm you
    analyzed in parts (a)--(c) is the most natural way to solve the
    problem---after all, it just iterates through the relevant entries of the
    array $B$, filling a value for each---it contains some highly unnecessary
    sources of inefficiency. Give a different algorithm to solve this problem,
    with an asymptotically better running time. 


\begin{Solution}
The trick here is to recognize that we can reuse computation from previous
steps. Rather than recompute entries $A[i]$ through $A[j-1]$, simply use what is
stored in $B[i,j-1]$ and add $A[j]$ to it.

Because we want the sum from $A[i]$ through $A[j]$ inclusive, we have to
initialize the first diagonal to $A[i] + A[i+1]$ separately.

\begin{algo}
   For $i = 1, 2, \ldots, n$\+ \\
   $B[i,i+1] = A[i] + A[i+1]$ \\
   For $j = i + 2, i + 3, \ldots, n$\+ \\
   $B[i,j] = B[i,j-1] + A[j]$\- \\
\end{algo}
\end{Solution}

\item What does ``asymptotically better'' mean?
\begin{Solution}
I didn't grade this problem, since I wasn't very clear about what I wanted. I hoped
you might think about how your first algorithm was $\Theta(n^3)$ but your second was $\Theta(n^2)$,
and notice that those functions were different, so you could say that an algorithm with runtime $f(n) = \Theta(g(n))$
is asymptotically better than an algorithm with runtime $f'(n) = \Theta(h(n))$ if $g(n) = \Omega(h(n))$ 
but $h(n) \neq \Omega(g(n))$, or something of this sort.

\end{Solution}

\end{enumerate}


\end{enumerate}

\end{document}

